\chapter{\label{ap:tor}Tor high performance computing cluster}
\graphicspath{{}}

\begin{figure}
    \begin{center}
        \input{Memory_usage_plot.pgf}
    \end{center}
    \caption{\label{fig:Memory_usage_plot}
    The memory storage of the matrix representation of the Hamiltonian Equation \ref{eq:spin_triplet_BdG_Hamiltonian} in \texttt{Python}'s \texttt{NumPy} library. Complex matrix entries are $16$ bytes and the size of the matrix depends on the number of sites $N=N_x=N_y$ as $\left(2^3\cdot N^2\right)^2$. Typical RAM capacity ($64$gb) limits us to $N=87$ ($58.7$gb).
    }
\end{figure}

The Tor cluster is a computing server located in the School of Physical Sciences at the University of Kent. It consists of 36 computational nodes, 32 of which provide 64 gigabytes of random-access memory (RAM). The model Hamiltonian described in Section \ref{sec:Non-unitary_spin-triplet} is represented by a matrix with $(50 \times 50 \text{ (sites)} \times 2 \text{ (orbitals)} \times 2 \text{ (spin)} \times 2 \text{ (particle-hole)})^2 = 4\times 10^8$ entries. The matrix was solved using the linear algebra Python library NumPy, which allocates 80 bytes to the Complex128 data type. Our simulations therefore require $3.2\times 10^{10} \text{ bytes} \sim 30 \text{ gigabytes}$ of RAM. Therefore we should be able to solve two model matrices simultaneously. However, due to overhead, we are only able to run one simulation at a time.

First a meta configuration file is fed to a meta Python script which produces configuration files. These configuration files are then fed to a Python script which creates a matrix representation of a model Hamiltonian and returns the solutions of the associated eigenproblem. Other Python scripts are then used to plot functions of the eigensystem, such as Green's functions, topography, spectrum, etc.

The meta configuration file is a .conf document written containing, for example, the following text:
\begin{lstlisting}[language=Python]
41,41,1 #n_sites_x min max step
41,41,1 #n_sites_y
2 #n_spins
2 #n_orbitals
True #closed_BC = False/True
-3.99,-2.87,0.21 #mu
0,0,1 #s
0,0,1 #delta
1.0,1.0,1 # t
0,0,1 # Delta
-4.4,4.4,1.1 # V
1,1,1 # n_impurities
0,0,0 #eta fixed
\end{lstlisting}
The meta Python script will produce configuration files from the \textit{min}imum to the \textit{max}imum values in given \textit{step} sizes.

A Bash script feeds the configuration files through the Python script which creates the associated Hamiltonian matrix and returns its eigensystem as a NumPy zipped data file named in the format \textit{n\_sites\_x, n\_sites\_y, mu, s, delta, t, Delta, V, n\_impurities}, e.g. `41\_41\_-2.94\_0\_0\_1.00\_0\_0\_1.conf'. The eigensystem can be loaded into a new Python script in the following way:
\begin{lstlisting}[language=Python]
eigensystem = np.load(`41_41_-2.94_0_0_1.00_0_0_1.conf')
eigenvalues = eigensystem['w']
eigenvectors  = eigensystem['v']
\end{lstlisting}
